\begin{proof}
The field $K$ is complete, nontrivially valued, and non-Archimedean by
construction.  The displayed Weierstrass equation has coefficients
$a=-1$, $b=0$, and discriminant
\[
 -16(4a^3+27b^2)=64\ne0.
\]
Thus $E_0$, with origin $[0:1:0]$, is an elliptic curve.  The constants
$\mathbf C\subset K$ lie in the valuation ring $K^\circ$, and
\[
 \mathcal A=E_0\times_{\mathbf C}\operatorname{Spec}K^\circ
\]
is a smooth proper group scheme over $K^\circ$, since smoothness and
properness are preserved by base change.  Its generic fibre is $A$, so
$A$ has good reduction.

For an abelian variety with good reduction, the torus in its Raynaud
uniformization is trivial; see \cite[\S16]{GJR26}.  Hence its cocharacter
lattice $N$ is zero and
\[
 M_{\mathbf R}=\operatorname{Hom}(N,\mathbf R)=0,
 \qquad
 \Lambda^{1,1}M_{\mathbf R}
 =\Lambda^1M_{\mathbf R}\otimes_{\mathbf R}\Lambda^1M_{\mathbf R}=0.
\]
On the other hand, \cite[Theorem~1.14]{GJR26} states that for a smooth
projective curve $Y$, the weakly smooth Dolbeault group
$H^{p,q}(Y^{\An})$ has real dimension one when $p+q\in\{0,2\}$.
Applying this with $Y=A$ and $(p,q)=(1,1)$ gives the second assertion in
\eqref{eq:ch18-gjr-cohomology-counterexample}.  The injection from
\cite[\S1.15]{GJR26} therefore has zero source and a one-dimensional target,
and cannot be surjective.

The toric rank of this example is zero, whereas a totally degenerate elliptic
curve has toric rank one.  Thus the argument answers the literal unrestricted
question but makes no assertion about the totally degenerate refinement.
\end{proof}

% PACKET-ONLY PRIMARY-SOURCE ATTESTATION; NOT LIVE CHAPTER TEXT
Primary source: Walter Gubler, Philipp Jell, and Joseph Rabinoff,
"Forms on Berkovich spaces based on harmonic tropicalizations," Bulletin de
la Societe Mathematique de France 154 (2026), 135--238,
DOI 10.24033/bsmf.2919; arXiv:2111.05741v3.

Theorem 1.14 (published paper, page 142 / PDF page 9): Let Y be a smooth
projective curve over K and let b be the first Betti number of X=Y^an.  Then
dim H^{p,q}(X) equals 1 if p+q is 0 or 2, and equals b if p+q=1.  These are
the Dolbeault groups of the weakly smooth complex introduced in the paper.

Section 1.15 (published paper, pages 142--143 / PDF pages 9--10): For an
abelian variety A, the canonical tropicalization uses the cocharacter lattice
N of the torus T in Raynaud uniformization.  With M_R=Hom(N,R) and
Lambda^{p,q}M_R=Lambda^p M_R tensor Lambda^q M_R, there is a canonical
injective homomorphism Lambda^{p,q}M_R -> H^{p,q}(A) for all p,q.  The authors
then write that they do not know whether this is an isomorphism as in the
complex case.

Raynaud-rank fact: if A has good reduction, the torus part T of its Raynaud
extension is trivial.  Thus N=0 and M_R=0.  A concrete example is obtained by
taking K to be the completed algebraic closure of C((t)) and taking A to be
the base change of any smooth elliptic curve over C; the constant smooth
proper model gives good reduction.

Scope warning: the statement in Section 1.15 is literally unrestricted.  A
counterexample with non-maximal degeneration answers that statement but does
not answer the refined question for totally degenerate abelian varieties.
